% sections/results.tex -- every measured value is a macro from numbers.tex. Test set unless stated otherwise.
\subsection*{Data audit}
FIVES contains \NDupPairsTrain{} pairs of training images and \NDupPairsTest{} pair of test images that
are byte-identical files with different masks (mask Dice over the whole image \DupTrainMaskDiceMin--\DupTrainMaskDiceMax{}
and \DupTestMaskDice), \NSharedMaskTestPairs{} further test pair that shares one mask between two near-identical images,
\NEmptyMasks{} empty training masks, \NSameEyePairs{} training image that shows the same eye as a test image (vessel-layout
IoU \SameEyeIoU; the files are not identical, and whether they are separate photographs cannot be determined from the
data), and \NBorderMasks{} masks with annotation lines outside the eye (Table~\ref{tab:data}, Fig~\ref{fig:audit}).
\NMarkings{} of the \NImages{} images carry burned-in text. \NIsoPairs{} isolated pairs of images share an identical marking
pattern and may come from one patient: \NIsoTrainTest{} pairs link our training set and \NIsoValTest{} our validation set
with the test set, \NIsoTestTest{} lie within the test set and none links training and validation. A further
\NAdjTestTest{} pairs of consecutively numbered test images share a marking pattern. \NNoMarking{} images carry no marking
and could not be checked in this way.

\subsection*{Overall segmentation}
The four CNNs reached nearly the same mean Dice, \CNNDiceMin--\CNNDiceMax{}, with standard deviations across seeds of
\CNNDiceSdSeedMin--\CNNDiceSdSeedMax{} (Table~\ref{tab:main}); at the image level their Dice differed from the U-Net's by
at most \MSDiceDiff{} (Table~\ref{tab:paired}), a negligible difference even where the Wilcoxon test detected it. The
fine-tuned SAM configurations had lower Dice (\SAMTiledDiceMean{} tiled, \SAMOverlapDiceMean{} with overlapping tiles,
\SAMWholeDiceMean{} whole image), and every tiled-SAM seed was below every U-Net seed in both Dice
(\SAMTiledDiceSeedMin--\SAMTiledDiceSeedMax{} against \UNetDiceSeedMin--\UNetDiceSeedMax) and NSD
(\SAMTiledNSDSeedMin--\SAMTiledNSDSeedMax{} against \UNetNSDSeedMin--\UNetNSDSeedMax). Tiled SAM also had lower precision
than the U-Net (\SAMTiledPrecMean{} against \UNetPrecMean; \nameref{si:overlap}). Image-level spread is given in
\nameref{si:imagelevel}.

% tables/main.tex -- generated by paper/scripts/make_assets.py. Do not edit by hand.
\begin{table}[!ht]
\begin{adjustwidth}{-0.5in}{-0.5in}
\centering
\caption{{\bf Test results.} Mean $\pm$ standard deviation across \NSeeds{} training seeds of each seed's mean over the \NTestImages{} test images, each run at its own threshold chosen on validation. Thin, medium and thick: centreline recall by ground-truth vessel radius ($<$ 3, 3--7, $>$ 7 px at 2048, 1-px tolerance). NSD: normalized surface Dice (boundary-pixel approximation, 2 px). Overlapping tiles: inference-only variant of the tiled SAM models.}
\begin{tabular}{lcccccc}
\hline
{\bf Model} & {\bf Dice} & {\bf clDice} & {\bf NSD} & {\bf Thin} & {\bf Medium} & {\bf Thick} \\ \thickhline
\multicolumn{7}{l}{\it CNNs (2 $\times$ 2 ablation)} \\ \hline
U-Net & \UNetDiceMean{} $\pm$ \UNetDiceSdSeed{} & \UNetClDiceMean{} $\pm$ \UNetClDiceSdSeed{} & \UNetNSDMean{} $\pm$ \UNetNSDSdSeed{} & \UNetThinMean{} $\pm$ \UNetThinSdSeed{} & \UNetMedMean{} $\pm$ \UNetMedSdSeed{} & \UNetThickMean{} $\pm$ \UNetThickSdSeed{} \\ \hline
Attention U-Net & \AttDiceMean{} $\pm$ \AttDiceSdSeed{} & \AttClDiceMean{} $\pm$ \AttClDiceSdSeed{} & \AttNSDMean{} $\pm$ \AttNSDSdSeed{} & \AttThinMean{} $\pm$ \AttThinSdSeed{} & \AttMedMean{} $\pm$ \AttMedSdSeed{} & \AttThickMean{} $\pm$ \AttThickSdSeed{} \\ \hline
Multi-Scale U-Net & \MSDiceMean{} $\pm$ \MSDiceSdSeed{} & \MSClDiceMean{} $\pm$ \MSClDiceSdSeed{} & \MSNSDMean{} $\pm$ \MSNSDSdSeed{} & \MSThinMean{} $\pm$ \MSThinSdSeed{} & \MSMedMean{} $\pm$ \MSMedSdSeed{} & \MSThickMean{} $\pm$ \MSThickSdSeed{} \\ \hline
Multi-Scale Attention U-Net & \MSAttDiceMean{} $\pm$ \MSAttDiceSdSeed{} & \MSAttClDiceMean{} $\pm$ \MSAttClDiceSdSeed{} & \MSAttNSDMean{} $\pm$ \MSAttNSDSdSeed{} & \MSAttThinMean{} $\pm$ \MSAttThinSdSeed{} & \MSAttMedMean{} $\pm$ \MSAttMedSdSeed{} & \MSAttThickMean{} $\pm$ \MSAttThickSdSeed{} \\ \hline
\multicolumn{7}{l}{\it Fine-tuned SAM (LoRA)} \\ \hline
SAM, tiled & \SAMTiledDiceMean{} $\pm$ \SAMTiledDiceSdSeed{} & \SAMTiledClDiceMean{} $\pm$ \SAMTiledClDiceSdSeed{} & \SAMTiledNSDMean{} $\pm$ \SAMTiledNSDSdSeed{} & \SAMTiledThinMean{} $\pm$ \SAMTiledThinSdSeed{} & \SAMTiledMedMean{} $\pm$ \SAMTiledMedSdSeed{} & \SAMTiledThickMean{} $\pm$ \SAMTiledThickSdSeed{} \\ \hline
SAM, overlapping tiles & \SAMOverlapDiceMean{} $\pm$ \SAMOverlapDiceSdSeed{} & \SAMOverlapClDiceMean{} $\pm$ \SAMOverlapClDiceSdSeed{} & \SAMOverlapNSDMean{} $\pm$ \SAMOverlapNSDSdSeed{} & \SAMOverlapThinMean{} $\pm$ \SAMOverlapThinSdSeed{} & \SAMOverlapMedMean{} $\pm$ \SAMOverlapMedSdSeed{} & \SAMOverlapThickMean{} $\pm$ \SAMOverlapThickSdSeed{} \\ \hline
SAM, whole image & \SAMWholeDiceMean{} $\pm$ \SAMWholeDiceSdSeed{} & \SAMWholeClDiceMean{} $\pm$ \SAMWholeClDiceSdSeed{} & \SAMWholeNSDMean{} $\pm$ \SAMWholeNSDSdSeed{} & \SAMWholeThinMean{} $\pm$ \SAMWholeThinSdSeed{} & \SAMWholeMedMean{} $\pm$ \SAMWholeMedSdSeed{} & \SAMWholeThickMean{} $\pm$ \SAMWholeThickSdSeed{} \\ \hline
\end{tabular}
\label{tab:main}
\end{adjustwidth}
\end{table}


% tables/paired.tex -- generated by paper/scripts/make_assets.py. Do not edit by hand.
\begin{table}[!ht]
\begin{adjustwidth}{-0.5in}{-0.5in}
\centering
\caption{{\bf Image-level comparison with U-Net.} Mean over the \NTestImages{} test images of the per-image difference from U-Net (each image's metric averaged over the model's seeds first), with the Holm-adjusted p-value of a two-sided paired Wilcoxon signed-rank test (correction within each metric across the six comparisons). These tests treat images as the sample and do not include training-seed variability.}
\begin{tabular}{lcccc}
\hline
{\bf Model $-$ U-Net} & {\bf Dice} & {\bf Thin} & {\bf clDice} & {\bf NSD} \\ \thickhline
Attention U-Net & \AttDiceDiff{} (\AttDiceP{}) & \AttThinDiff{} (\AttThinP{}) & \AttClDiceDiff{} (\AttClDiceP{}) & \AttNSDDiff{} (\AttNSDP{}) \\ \hline
Multi-Scale U-Net & \MSDiceDiff{} (\MSDiceP{}) & \MSThinDiff{} (\MSThinP{}) & \MSClDiceDiff{} (\MSClDiceP{}) & \MSNSDDiff{} (\MSNSDP{}) \\ \hline
Multi-Scale Attention U-Net & \MSAttDiceDiff{} (\MSAttDiceP{}) & \MSAttThinDiff{} (\MSAttThinP{}) & \MSAttClDiceDiff{} (\MSAttClDiceP{}) & \MSAttNSDDiff{} (\MSAttNSDP{}) \\ \hline
SAM, tiled & \SAMTiledDiceDiff{} (\SAMTiledDiceP{}) & \SAMTiledThinDiff{} (\SAMTiledThinP{}) & \SAMTiledClDiceDiff{} (\SAMTiledClDiceP{}) & \SAMTiledNSDDiff{} (\SAMTiledNSDP{}) \\ \hline
SAM, overlapping tiles & \SAMOverlapDiceDiff{} (\SAMOverlapDiceP{}) & \SAMOverlapThinDiff{} (\SAMOverlapThinP{}) & \SAMOverlapClDiceDiff{} (\SAMOverlapClDiceP{}) & \SAMOverlapNSDDiff{} (\SAMOverlapNSDP{}) \\ \hline
SAM, whole image & \SAMWholeDiceDiff{} (\SAMWholeDiceP{}) & \SAMWholeThinDiff{} (\SAMWholeThinP{}) & \SAMWholeClDiceDiff{} (\SAMWholeClDiceP{}) & \SAMWholeNSDDiff{} (\SAMWholeNSDP{}) \\ \hline
\end{tabular}
\label{tab:paired}
\end{adjustwidth}
\end{table}


\subsection*{Thin vessels}
Thin vessels were the hardest class for every model, with recall from \AllThinMin{} to \AllThinMax{} across the seven
configurations against \AllMedMin--\AllMedMax{} for medium and \AllThickMinP--\AllThickMaxP{} for thick vessels
(Table~\ref{tab:main}). The ground truth pushed through the 512 $\times$ 512 grid and back kept \RefFiveThin{} of the
thin-vessel centreline, and through the 256 $\times$ 256 grid of whole-image SAM \RefTwoThin{} (\nameref{si:oracle}). With
cut-offs of 4 and 8 pixels the ordering of the models was unchanged, but the differences shrank: the tiled-SAM advantage
over the U-Net fell from \SAMTiledThinDiff{} to \SAMTiledThinAltDiff{} (\nameref{si:width}), so it lies mainly in the
thinnest vessels.

\subsection*{Attention gates and multi-scale blocks}
Neither module changed Dice by more than \AblDiceMainMaxAbs{} (main effects, Table~\ref{tab:ablation},
Fig~\ref{fig:ablation}). The attention main effect on thin-vessel recall changed sign with the threshold rule: \AblAttThinSel{}
(95\% CI \AblAttThinSelLo{} to \AblAttThinSelHi) at the validation-selected thresholds and \AblAttThinFix{}
(\AblAttThinFixLo{} to \AblAttThinFixHi) at 0.5, with precision and recall moving in opposite directions under the two
rules, and its models' selected thresholds were equal or lower (\nameref{si:runs}). Attention therefore had no
consistent effect on thin-vessel recall.

Multi-scale blocks lowered thin-vessel recall under both rules: \AblMSThinSel{} (\AblMSThinSelLo{} to \AblMSThinSelHi) and
\AblMSThinFix{} (\AblMSThinFixLo{} to \AblMSThinFixHi). The decrease was consistent in direction but overlapped
training-seed variability (Fig~\ref{fig:ablation}D, E): per-seed thin-vessel recall was \MSThinSeedMin--\MSThinSeedMax{} for
the Multi-Scale U-Net and \UNetThinSeedMin--\UNetThinSeedMax{} for the U-Net, and at 0.5 the best Multi-Scale U-Net seed
(\MSThinFixSeedMax) exceeded every U-Net seed. The modules also interacted. At 0.5 the thin-vessel interaction
(\AblIntThinFix, \AblIntThinFixLo{} to \AblIntThinFixHi) was as large as either main effect: attention changed thin-vessel
recall by \SimpleAttNoMSThinFix{} without multi-scale blocks and by \SimpleAttWithMSThinFix{} with them. At the selected
thresholds the precision interaction (\AblIntPrecSel) exceeded both main effects (\nameref{si:simple}).

% tables/ablation.tex -- generated by paper/scripts/make_assets.py. Do not edit by hand.
\begin{table}[!ht]
\begin{adjustwidth}{-0.5in}{-0.5in}
\centering
\caption{{\bf 2 $\times$ 2 ablation of attention gates and multi-scale blocks.} Effects on test metrics with 95\% bootstrap confidence intervals (resampling the \NTestImages{} images; each image averaged over seeds). Attention main effect: mean of (Attention U-Net $-$ U-Net) and (Multi-Scale Attention U-Net $-$ Multi-Scale U-Net); multi-scale main effect likewise; interaction: (Multi-Scale Attention U-Net $-$ Multi-Scale U-Net) $-$ (Attention U-Net $-$ U-Net). clDice and NSD exist only at the val-selected thresholds.}
\begin{tabular}{lccc}
\hline
{\bf Metric} & {\bf Attention} & {\bf Multi-scale} & {\bf Interaction} \\ \thickhline
\multicolumn{4}{l}{\it Threshold: val-selected} \\ \hline
Dice & \AblAttDiceSel{} [\AblAttDiceSelLo{}, \AblAttDiceSelHi{}] & \AblMSDiceSel{} [\AblMSDiceSelLo{}, \AblMSDiceSelHi{}] & \AblIntDiceSel{} [\AblIntDiceSelLo{}, \AblIntDiceSelHi{}] \\ \hline
Thin recall & \AblAttThinSel{} [\AblAttThinSelLo{}, \AblAttThinSelHi{}] & \AblMSThinSel{} [\AblMSThinSelLo{}, \AblMSThinSelHi{}] & \AblIntThinSel{} [\AblIntThinSelLo{}, \AblIntThinSelHi{}] \\ \hline
clDice & \AblAttClDiceSel{} [\AblAttClDiceSelLo{}, \AblAttClDiceSelHi{}] & \AblMSClDiceSel{} [\AblMSClDiceSelLo{}, \AblMSClDiceSelHi{}] & \AblIntClDiceSel{} [\AblIntClDiceSelLo{}, \AblIntClDiceSelHi{}] \\ \hline
NSD & \AblAttNSDSel{} [\AblAttNSDSelLo{}, \AblAttNSDSelHi{}] & \AblMSNSDSel{} [\AblMSNSDSelLo{}, \AblMSNSDSelHi{}] & \AblIntNSDSel{} [\AblIntNSDSelLo{}, \AblIntNSDSelHi{}] \\ \hline
Precision & \AblAttPrecSel{} [\AblAttPrecSelLo{}, \AblAttPrecSelHi{}] & \AblMSPrecSel{} [\AblMSPrecSelLo{}, \AblMSPrecSelHi{}] & \AblIntPrecSel{} [\AblIntPrecSelLo{}, \AblIntPrecSelHi{}] \\ \hline
Recall & \AblAttRecSel{} [\AblAttRecSelLo{}, \AblAttRecSelHi{}] & \AblMSRecSel{} [\AblMSRecSelLo{}, \AblMSRecSelHi{}] & \AblIntRecSel{} [\AblIntRecSelLo{}, \AblIntRecSelHi{}] \\ \hline
\multicolumn{4}{l}{\it Threshold: fixed 0.5} \\ \hline
Dice & \AblAttDiceFix{} [\AblAttDiceFixLo{}, \AblAttDiceFixHi{}] & \AblMSDiceFix{} [\AblMSDiceFixLo{}, \AblMSDiceFixHi{}] & \AblIntDiceFix{} [\AblIntDiceFixLo{}, \AblIntDiceFixHi{}] \\ \hline
Thin recall & \AblAttThinFix{} [\AblAttThinFixLo{}, \AblAttThinFixHi{}] & \AblMSThinFix{} [\AblMSThinFixLo{}, \AblMSThinFixHi{}] & \AblIntThinFix{} [\AblIntThinFixLo{}, \AblIntThinFixHi{}] \\ \hline
Precision & \AblAttPrecFix{} [\AblAttPrecFixLo{}, \AblAttPrecFixHi{}] & \AblMSPrecFix{} [\AblMSPrecFixLo{}, \AblMSPrecFixHi{}] & \AblIntPrecFix{} [\AblIntPrecFixLo{}, \AblIntPrecFixHi{}] \\ \hline
Recall & \AblAttRecFix{} [\AblAttRecFixLo{}, \AblAttRecFixHi{}] & \AblMSRecFix{} [\AblMSRecFixLo{}, \AblMSRecFixHi{}] & \AblIntRecFix{} [\AblIntRecFixLo{}, \AblIntRecFixHi{}] \\ \hline
\end{tabular}
\label{tab:ablation}
\end{adjustwidth}
\end{table}


\begin{figure}[!h]
\figpreview{figures/fig_ablation.pdf}
\caption{\CapAblation}
\label{fig:ablation}
\end{figure}

\subsection*{Fine-tuned SAM}
Whole-image SAM had lower thin-vessel recall than the U-Net in every seed (\SAMWholeThinSeedMin--\SAMWholeThinSeedMax{}
against \UNetThinSeedMin--\UNetThinSeedMax; mean per-image difference \SAMWholeThinDiff; Fig~\ref{fig:sam}), and also lower
medium-vessel recall. Its thin-vessel recall exceeded that of its 256 $\times$ 256 reference.

Tiled SAM had the highest mean thin-vessel recall, \SAMTiledThinMean{} $\pm$ \SAMTiledThinSdSeed{} against \UNetThinMean{}
$\pm$ \UNetThinSdSeed{} for the U-Net (mean per-image difference \SAMTiledThinDiff, higher on \SAMTiledThinShare\% of
images). The gain lay within training-seed variability: the tiled-SAM seeds reached \SAMTiledThinSeedA, \SAMTiledThinSeedB{}
and \SAMTiledThinSeedC{} and the U-Net seeds \UNetThinSeedA, \UNetThinSeedB{} and \UNetThinSeedC, so one tiled-SAM seed was
below two U-Net seeds. Its clDice did not differ consistently from the U-Net's.

Across all three seeds the two SAM configurations differed consistently, in opposite directions on different metrics:
tiled SAM had higher thin-vessel recall, clDice and NSD and lower Dice than whole-image SAM (Fig~\ref{fig:sam}). Two tiled
runs and all three whole-image runs reached the 150-epoch limit with their best epoch near the end (\nameref{si:runs},
\nameref{si:training}).

\begin{figure}[!h]
\figpreview{figures/fig_sam.pdf}
\caption{\CapSam}
\label{fig:sam}
\end{figure}

On validation (seed 0), the non-overlapping tiles left a discontinuity at the tile borders (Dice difference within
$\pm$16~pixels \SeamSAMTiledDiDDice, \SeamSAMTiledDiDDiceLo{} to \SeamSAMTiledDiDDiceHi) that overlapping tiles removed
(\nameref{si:seam}). On the test set overlapping tiles raised Dice in every seed, though not to the U-Net's level, and left
thin-vessel recall unchanged at equal thresholds (\CurveOverlapMinusTiledMinTest{} to \CurveOverlapMinusTiledMaxTest).

\subsection*{Threshold sensitivity}
Thin-vessel recall depended strongly on the threshold (Fig~\ref{fig:threshold}); for the U-Net it fell from
\CurveUNetThinLow{} at 0.05 to \CurveUNetThinMid{} at 0.5 and \CurveUNetThinHigh{} at 0.95. The mean of the tiled-SAM seeds
lay above the best CNN at all \NThresholds{} thresholds (by \CurveSAMTiledGapBestMinTest{} to
\CurveSAMTiledGapBestMaxTest), but the range of its seeds overlapped that of the U-Net's seeds at
\CurveSAMTiledOverlapUNetTest{} of them; whole-image SAM lay below the CNNs except at the \CurveSAMWholeAboveBestTest{}
highest thresholds. Part of the U-Net's spread between seeds came from the threshold choice: its seed standard deviation
in thin-vessel recall was \UNetThinSdSeed{} at the selected thresholds and \UNetThinFixSdSeed{} at 0.5 (tiled SAM
\SAMTiledThinSdSeed{} and \SAMTiledThinFixSdSeed; \nameref{si:fixed}, \nameref{si:thrval}).

\begin{figure}[!h]
\figpreview{figures/fig_threshold.pdf}
\caption{\CapThreshold}
\label{fig:threshold}
\end{figure}

\subsection*{What Dice shows}
Mean Dice spanned \SpreadDice{} across the seven configurations and thin-vessel recall \SpreadThin{} (largest minus
smallest model mean; Fig~\ref{fig:tradeoff}A, B), similar spans relative to their seed variability, and Dice separated
every SAM configuration from the U-Net in every seed. The metrics ranked differently: whole-image SAM had higher Dice than
tiled SAM but much lower thin-vessel recall, and glaucoma images had the lowest Dice for every model but the highest
thin-vessel recall for the CNNs and whole-image SAM (\nameref{si:disease}).

Across seeds, the checkpoint-selection criterion did not determine thin-vessel recall (Fig~\ref{fig:tradeoff}C, D). The
three tiled-SAM runs reached a best validation Dice of \SAMTiledValBestDiceMin--\SAMTiledValBestDiceMax, while their
validation thin-vessel recall ranged from \SAMTiledValThinMin{} to \SAMTiledValThinMax{} at the selected thresholds and from
\SAMTiledValThinFixMin{} to \SAMTiledValThinFixMax{} at 0.5.

\begin{figure}[!h]
\figpreview{figures/fig_tradeoff.pdf}
\caption{\CapTradeoff}
\label{fig:tradeoff}
\end{figure}

\subsection*{Agreement between two available annotations}
For the \AgreeNPairs{} pairs of byte-identical images with two masks, the masks agreed closely by Dice inside the FOV
(\AgreeDiceMin--\AgreeDiceMax) but less on thin vessels: each mask recovered on average \AgreeThinMean{} of the other's
thin-vessel centreline (pairs \AgreeThinMin--\AgreeThinMax), against \AgreeMedMean{} for medium and \AgreeThickMean{} for
thick vessels (\nameref{si:agreement}). All pairs are AMD images, and n = \AgreeNPairs.

\subsection*{Qualitative examples and secondary results}
In native-resolution crops (Fig~\ref{fig:qualitative}, \nameref{si:qualall}), false negatives were mostly thin branches and
false positives mostly narrow rims along larger vessels. Normal images had the lowest thin-vessel recall for every model (\nameref{si:disease}, \nameref{si:diseasefig}). All models slightly underestimated vessel
density and fractal dimension; the largest fractal-dimension errors occurred on two glaucoma images with very sparse
annotation (vessel density \FDOutOneDens{} and \FDOutTwoDens; \nameref{si:vessel}, \nameref{si:blandaltman}). Tiled SAM
needed \SAMTiledInfMs{}~ms per image, \SAMTiledInfRatio{} times the U-Net (Table~\ref{tab:cost}, \nameref{si:costfig}).
Zero-shot SAM with oracle prompts labelled most of the FOV as vessel (Dice \OracleDice; \nameref{si:oracle}).

\begin{figure}[!h]
\figpreview{figures/fig_qualitative_raster.png}
\caption{\CapQualitative}
\label{fig:qualitative}
\end{figure}

% tables/cost.tex -- generated by paper/scripts/make_assets.py. Do not edit by hand.
\begin{table}[!ht]
\begin{adjustwidth}{-0.5in}{-0.5in}
\centering
\caption{{\bf Computational cost.} One NVIDIA L40S GPU. Inference: one image per call over the \NTestImages{} test images after 10 warm-up calls, from the input on the GPU to the probability map (tiling, stitching or blending included); FLOPs per image from PyTorch's FlopCounterMode; peak memory allocated. Training time and memory: seed-0 runs.}
\begin{tabular}{lcccccc}
\hline
{\bf Model} & {\bf Trainable params (M)} & {\bf Inference (ms/image)} & {\bf GFLOPs/image} & {\bf Inference memory (GB)} & {\bf Training time (h)} & {\bf Training memory (GB)} \\ \thickhline
U-Net & \UNetParamsTrainableM{} & \UNetInfMs{} & \UNetGflops{} & \UNetInfMem{} & \UNetTrainH{} & \UNetTrainMem{} \\ \hline
Attention U-Net & \AttParamsTrainableM{} & \AttInfMs{} & \AttGflops{} & \AttInfMem{} & \AttTrainH{} & \AttTrainMem{} \\ \hline
Multi-Scale U-Net & \MSParamsTrainableM{} & \MSInfMs{} & \MSGflops{} & \MSInfMem{} & \MSTrainH{} & \MSTrainMem{} \\ \hline
Multi-Scale Attention U-Net & \MSAttParamsTrainableM{} & \MSAttInfMs{} & \MSAttGflops{} & \MSAttInfMem{} & \MSAttTrainH{} & \MSAttTrainMem{} \\ \hline
SAM, tiled & \SAMTiledParamsTrainableM{} & \SAMTiledInfMs{} & \SAMTiledGflops{} & \SAMTiledInfMem{} & \SAMTiledTrainH{} & \SAMTiledTrainMem{} \\ \hline
SAM, overlapping tiles & \SAMOverlapParamsTrainableM{} & \SAMOverlapInfMs{} & \SAMOverlapGflops{} & \SAMOverlapInfMem{} & (tiled models) & -- \\ \hline
SAM, whole image & \SAMWholeParamsTrainableM{} & \SAMWholeInfMs{} & \SAMWholeGflops{} & \SAMWholeInfMem{} & \SAMWholeTrainH{} & \SAMWholeTrainMem{} \\ \hline
\end{tabular}
\label{tab:cost}
\end{adjustwidth}
\end{table}


\subsection*{Post hoc analyses}
At matched precision the tiled-SAM advantage in thin-vessel recall shrank from \SAMTiledThinDiff{} to
\PHStTiledThinMatch{} (\nameref{si:matched}, \nameref{si:precthin}), because tiled SAM had operated at lower precision. Whole-image SAM stayed lower (\PHStWholeThinMatch) and the multi-scale decrease persisted
(\PHStMSThinMatch). Thin-vessel precision was \PHWPrecThinMin--\PHWPrecThinMax{} for all models, against
\PHWPrecMedMin--\PHWPrecMedMax{} for medium and \PHWPrecThickMin--\PHWPrecThickMax{} for thick vessels; multi-scale blocks
raised it (\PHStMSWPrecSel) while lowering thin-vessel recall. At 0.5, tiled SAM had lower clDice (\PHStTiledClDiceFix)
and NSD (\PHStTiledNSDFix) than the U-Net in every seed (\nameref{si:precision}). Pruning skeleton spurs removed
\PHPrunedThinShare\% of the thin-vessel centreline and barely changed the comparisons (tiled SAM \PHStTiledThinPrSel,
multi-scale effect \PHStMSThinPrSel).

Resampling \PHNClusters{} clusters of possibly same-patient test images instead of images barely changed the intervals;
also resampling seeds widened them, for thin-vessel recall several-fold (\nameref{si:stats}). These
intervals excluded zero for the Dice and NSD deficits of tiled SAM, the thin-vessel deficit of whole-image SAM and the
multi-scale decrease in thin-vessel recall, which was also negative in all three coupled seed pairs under every threshold
rule, but included zero for the tiled-SAM thin-vessel gain and every attention effect on thin-vessel recall and
Dice.
